\section{Conclusion}\label{conclusion}

This paper presented a reusable, configuration-driven local benchmarking pipeline for corpus-specific embedding-model selection in RAG retrieval. The workflow takes an unlabeled local document collection as its starting point, constructs the queries and relevance judgments needed for evaluation, and produces the rankings, uncertainty estimates, and reports required for corpus-specific candidate comparison. Each collection is processed separately through ingestion, chunking, synthetic query generation, pooled qrels construction, multi-model retrieval evaluation, timing capture, and automated reporting. Its central research position is deliberately bounded: practitioners should not assume that one embedding model will be optimal for every corpus under their control. The three-corpus snapshot demonstrates how the workflow tests that assumption; it does not prove that a universal best model cannot exist.

For RQ1, the fixed snapshot shows that the configurable pipeline can produce inspectable, corpus-specific decision-support evidence. Separate reports for \texttt{education\_eng}, \texttt{education\_tr}, and \texttt{law\_tr} contain comparable retrieval metrics, timing summaries, query-level uncertainty, and preserved artifacts under the same five-model protocol. These reports support corpus-specific screening rather than a single cross-corpus selection rule.

For RQ2, the two education cases show language-associated ranking differences: English-oriented candidates lead \texttt{education\_eng}, whereas multilingual candidates lead \texttt{education\_tr}. This comparison does not establish a causal language effect because each language condition is represented by one corpus. For RQ3, the two Turkish cases instead show domain-associated stability: \texttt{M5} remains the quality-first candidate and \texttt{M3} the lower-latency alternative for both \texttt{education\_tr} and \texttt{law\_tr}. This stability does not show that domain is irrelevant; together, RQ2 and RQ3 illustrate why both shortlist change and shortlist stability should be examined through separate corpus runs.

For RQ4, the reports show that retrieval quality and operational priorities can produce different corpus-specific selections. They combine retrieval metrics with compound corpus embedding/index-insertion throughput, compound query-embedding/vector-search latency, and uncertainty intervals. What the uncertainty layer adds is visible in the demonstrated scenario: on \texttt{education\_eng} the quality advantage of the leading candidate over its runner-up lies within the paired bootstrap interval, so latency becomes the practical differentiator, whereas on both Turkish corpora the corresponding advantage is separated and the observed quality evidence guides the selection. The aggregate performance of \texttt{M5} remains secondary comparative context and does not override these corpus-level readings.

The current limitations also define clear directions for further evaluation. Human audit and judge calibration would strengthen the validation of synthetic queries and LLM-generated relevance judgments; controlled sensitivity studies would test dependence on retrieval and benchmark configuration; broader corpus and model coverage would assess transfer beyond the present snapshot; and repeated operational measurements would support stronger deployment-oriented performance claims.

The practical takeaway is that practitioners should not assume that a model selected from public benchmarks or a previous corpus will transfer unchanged to a new collection; materially different corpora should be evaluated using corpus-specific evidence, with aggregate results treated as secondary context and the final selection guided by project priorities.
